% sec-01: the independent-scientist case, ARPA framing.  Sources:
% science-golden-age chunk-01 and chunk-03 (ARPA program-manager model).
\section{One Performer, Bounded Milestones}

\textit{Science: A New Golden Age} asks the federal research portfolio, roughly
\textbf{\$200 billion} a year, to stop deferring to ``the same incumbents that
consume the funding'' and to invest directly in individual researchers
\cite{kratsios2026sciencegoldenage}. Chapter II names the instrument it trusts
for high-risk work: the ARPA program-manager model, with ARPA-H and ARPA-E as
the live examples, chosen because a program manager can start a bet, run it on
milestones, and end it without a committee.

This proposal is written to be run that way. It is one performer, not a
consortium. Every milestone below has a stated kill criterion, and each is
checkable from data the program itself produces rather than from a narrative
progress report. The work that would normally consume the first eighteen months
of an ARPA program, the protocol, the regulatory package, and the modeling, is
already finished and published \cite{foundingdocs}.

\begin{apptable}
\begin{tabularx}{\textwidth}{>{\raggedright\arraybackslash}p{4.35cm} >{\raggedright\arraybackslash}X}
\headrow \thead{ARPA practice the report endorses} & \thead{How this program is written for it}\\
\midrule
A program manager who can end a bet & Three gates, each with a numeric kill criterion, not a narrative review \\
Milestones tied to measurable outputs & Stop latency in milliseconds, conversion counts, dose-limiting toxicity counts \\
Performers outside the incumbent set & One independent scientist contracting a single clinical site \\
A defined transition out of the program & An RP2D, a public safety set, and an open verification package any site can reuse \\
\end{tabularx}
\end{apptable}
\tabcap{Four features of the ARPA model that the report endorses, and the
specific way each is written into this program so that it can be checked at a
gate rather than argued at a review.}
